\documentclass[11pt,a4paper]{article}
\newcommand{\intestazione}{Dimero di spin --- 8. Correlatori sotto rumore}
% =====================================================================
%  Preambolo comune ai documenti sul dimero (serie "che cosa abbiamo fatto")
%  Incluso con % =====================================================================
%  Preambolo comune ai documenti sul dimero (serie "che cosa abbiamo fatto")
%  Incluso con % =====================================================================
%  Preambolo comune ai documenti sul dimero (serie "che cosa abbiamo fatto")
%  Incluso con \input{preambolo_dimero} subito dopo \documentclass.
% =====================================================================
\usepackage[T1]{fontenc}
\usepackage[utf8]{inputenc}
\usepackage[main=italian,provide=*]{babel}
\usepackage{amsmath,amssymb,mathtools}
\usepackage{graphicx}
\usepackage{booktabs}
\usepackage{colortbl}
\usepackage{xcolor}
\usepackage{geometry}
\usepackage{placeins}
\usepackage{caption}
\usepackage{enumitem}
\usepackage{fancyhdr}
\usepackage{needspace}
\usepackage{tikz}
\usetikzlibrary{arrows.meta,positioning,shapes.geometric,calc,fit,backgrounds}
\usepackage{hyperref}
\hypersetup{colorlinks=true,linkcolor=blu,citecolor=blu,urlcolor=blu}

\geometry{margin=2.5cm}

% --- colori coerenti con quelli delle figure (palette Okabe-Ito)
\definecolor{blu}{HTML}{0072B2}
\definecolor{arancio}{HTML}{E69F00}
\definecolor{verdes}{HTML}{009E73}
\definecolor{rossos}{HTML}{D55E00}
\definecolor{violas}{HTML}{CC79A7}
\definecolor{grigetto}{HTML}{E8E8E8}

\captionsetup{font=small,labelfont=bf,width=0.92\textwidth}

\pagestyle{fancy}
\fancyhf{}
\fancyhead[L]{\small\itshape\intestazione}
\fancyhead[R]{\small\thepage}
\renewcommand{\headrulewidth}{0.3pt}

% --- notazione
\newcommand{\ket}[1]{\lvert #1\rangle}
\newcommand{\bra}[1]{\langle #1 \rvert}
\newcommand{\media}[1]{\langle #1 \rangle}
\newcommand{\vs}[1]{\mathbf{s}_{#1}}

% --- riquadro "in una riga": il messaggio da portare a casa
\newcommand{\messaggio}[1]{%
  \begin{center}
  \begin{minipage}{0.93\textwidth}
  \setlength{\fboxsep}{9pt}%
  \colorbox{grigetto}{\begin{minipage}{\dimexpr\textwidth-2\fboxsep}
  \small\textbf{In una riga.}\enspace #1
  \end{minipage}}
  \end{minipage}
  \end{center}
}

% --- riquadro "come lo sappiamo": la verifica numerica dietro un'affermazione
\newcommand{\verifica}[1]{%
  \begin{center}
  \begin{minipage}{0.93\textwidth}
  \small\textcolor{verdes}{\rule{2.2em}{1.1pt}}\enspace
  \textcolor{verdes}{\textbf{Come lo sappiamo.}}\enspace #1
  \end{minipage}
  \end{center}
  \medskip
}
 subito dopo \documentclass.
% =====================================================================
\usepackage[T1]{fontenc}
\usepackage[utf8]{inputenc}
\usepackage[main=italian,provide=*]{babel}
\usepackage{amsmath,amssymb,mathtools}
\usepackage{graphicx}
\usepackage{booktabs}
\usepackage{colortbl}
\usepackage{xcolor}
\usepackage{geometry}
\usepackage{placeins}
\usepackage{caption}
\usepackage{enumitem}
\usepackage{fancyhdr}
\usepackage{needspace}
\usepackage{tikz}
\usetikzlibrary{arrows.meta,positioning,shapes.geometric,calc,fit,backgrounds}
\usepackage{hyperref}
\hypersetup{colorlinks=true,linkcolor=blu,citecolor=blu,urlcolor=blu}

\geometry{margin=2.5cm}

% --- colori coerenti con quelli delle figure (palette Okabe-Ito)
\definecolor{blu}{HTML}{0072B2}
\definecolor{arancio}{HTML}{E69F00}
\definecolor{verdes}{HTML}{009E73}
\definecolor{rossos}{HTML}{D55E00}
\definecolor{violas}{HTML}{CC79A7}
\definecolor{grigetto}{HTML}{E8E8E8}

\captionsetup{font=small,labelfont=bf,width=0.92\textwidth}

\pagestyle{fancy}
\fancyhf{}
\fancyhead[L]{\small\itshape\intestazione}
\fancyhead[R]{\small\thepage}
\renewcommand{\headrulewidth}{0.3pt}

% --- notazione
\newcommand{\ket}[1]{\lvert #1\rangle}
\newcommand{\bra}[1]{\langle #1 \rvert}
\newcommand{\media}[1]{\langle #1 \rangle}
\newcommand{\vs}[1]{\mathbf{s}_{#1}}

% --- riquadro "in una riga": il messaggio da portare a casa
\newcommand{\messaggio}[1]{%
  \begin{center}
  \begin{minipage}{0.93\textwidth}
  \setlength{\fboxsep}{9pt}%
  \colorbox{grigetto}{\begin{minipage}{\dimexpr\textwidth-2\fboxsep}
  \small\textbf{In una riga.}\enspace #1
  \end{minipage}}
  \end{minipage}
  \end{center}
}

% --- riquadro "come lo sappiamo": la verifica numerica dietro un'affermazione
\newcommand{\verifica}[1]{%
  \begin{center}
  \begin{minipage}{0.93\textwidth}
  \small\textcolor{verdes}{\rule{2.2em}{1.1pt}}\enspace
  \textcolor{verdes}{\textbf{Come lo sappiamo.}}\enspace #1
  \end{minipage}
  \end{center}
  \medskip
}
 subito dopo \documentclass.
% =====================================================================
\usepackage[T1]{fontenc}
\usepackage[utf8]{inputenc}
\usepackage[main=italian,provide=*]{babel}
\usepackage{amsmath,amssymb,mathtools}
\usepackage{graphicx}
\usepackage{booktabs}
\usepackage{colortbl}
\usepackage{xcolor}
\usepackage{geometry}
\usepackage{placeins}
\usepackage{caption}
\usepackage{enumitem}
\usepackage{fancyhdr}
\usepackage{needspace}
\usepackage{tikz}
\usetikzlibrary{arrows.meta,positioning,shapes.geometric,calc,fit,backgrounds}
\usepackage{hyperref}
\hypersetup{colorlinks=true,linkcolor=blu,citecolor=blu,urlcolor=blu}

\geometry{margin=2.5cm}

% --- colori coerenti con quelli delle figure (palette Okabe-Ito)
\definecolor{blu}{HTML}{0072B2}
\definecolor{arancio}{HTML}{E69F00}
\definecolor{verdes}{HTML}{009E73}
\definecolor{rossos}{HTML}{D55E00}
\definecolor{violas}{HTML}{CC79A7}
\definecolor{grigetto}{HTML}{E8E8E8}

\captionsetup{font=small,labelfont=bf,width=0.92\textwidth}

\pagestyle{fancy}
\fancyhf{}
\fancyhead[L]{\small\itshape\intestazione}
\fancyhead[R]{\small\thepage}
\renewcommand{\headrulewidth}{0.3pt}

% --- notazione
\newcommand{\ket}[1]{\lvert #1\rangle}
\newcommand{\bra}[1]{\langle #1 \rvert}
\newcommand{\media}[1]{\langle #1 \rangle}
\newcommand{\vs}[1]{\mathbf{s}_{#1}}

% --- riquadro "in una riga": il messaggio da portare a casa
\newcommand{\messaggio}[1]{%
  \begin{center}
  \begin{minipage}{0.93\textwidth}
  \setlength{\fboxsep}{9pt}%
  \colorbox{grigetto}{\begin{minipage}{\dimexpr\textwidth-2\fboxsep}
  \small\textbf{In una riga.}\enspace #1
  \end{minipage}}
  \end{minipage}
  \end{center}
}

% --- riquadro "come lo sappiamo": la verifica numerica dietro un'affermazione
\newcommand{\verifica}[1]{%
  \begin{center}
  \begin{minipage}{0.93\textwidth}
  \small\textcolor{verdes}{\rule{2.2em}{1.1pt}}\enspace
  \textcolor{verdes}{\textbf{Come lo sappiamo.}}\enspace #1
  \end{minipage}
  \end{center}
  \medskip
}


\title{\bfseries Quando il rumore incontra una misura vera\\[2pt]
\large Documento 8 --- correlatori dinamici sotto rumore}
\author{Samuele Schiavi \\ \small Tesi triennale in Fisica, Università di Parma
--- Relatore: Prof.\ Alessandro Chiesa}
\date{}

\begin{document}
\maketitle
\thispagestyle{fancy}

\section{Una differenza qualitativa rispetto ai due documenti precedenti}

Il VQE rumoroso (Documento 6) e la fedeltà di Trotter rumorosa (Documento 7)
sono entrambi oggetti che si leggono direttamente dall'operatore densità:
un'aspettazione, una fedeltà. Un correlatore dinamico
$C_{ij}^{\alpha\beta}(t)=\bra{\psi_0}\sigma_i^\alpha(t)\sigma_j^\beta(0)\ket{\psi_0}$
non è così: come già nel Documento 4, si estrae da un circuito di Hadamard
test con un qubit ausiliario, e la quantità finale passa da una \emph{misura
reale} sull'ancilla. È il primo punto della pipeline in cui il readout, fin
qui rimasto sullo sfondo del modello di rumore (Documento 5), entra
effettivamente in gioco.

\section{Il circuito, con il rumore su ogni ingrediente}

Lo schema è lo stesso del Documento 4: preparazione di $\ket{\psi_0}$,
Hadamard sull'ancilla, controlled-$W$ sul sito $j$ prima di $U(t)$, evoluzione
$U(t)$ non controllata (Trotter, come nel Documento 7), anti-controlled-$V$ sul
sito $i$ dopo $U(t)$, rotazione di base e misura $Z$ sull'ancilla. Sotto
rumore, ciascuno di questi ingredienti è imperfetto: la preparazione (Documento
6), il blocco $U(t)$ ripetuto (Documento 7), e infine la misura, che include
anche l'errore di lettura.

La preparazione può usare le ampiezze esatte oppure il circuito VQE reale ---
la stessa opzione già in uso per anello e catena a $N=3$. Verificarla di nuovo
in questa sessione si è rivelato necessario, non solo prudente: vedi
Sez.~\ref{sec:regressione}.

\section{Il readout: correzione simmetrica}

\begin{equation}
\langle Z\rangle_\text{readout} = (1-2p)\,\langle Z\rangle_\text{ideale}
\end{equation}
esattamente come nel Documento 5 (il caso asimmetrico è un'estensione a sé,
Documento 11). Un fattore puramente moltiplicativo: la ragione, resa esplicita
al documento successivo, per cui $N^*$ non dipenderà da $p_\text{readout}$.

\section{Una regressione trovata e richiusa}
\label{sec:regressione}

Il modulo che implementa la preparazione via ansatz VQE reale, usato per il
confronto di coerenza qui sotto, era regredito a una versione precedente ---
un placeholder senza implementazione al posto del codice giusto. Non un errore
mai corretto: un file caricato per sbaglio sopra la versione funzionante, un
episodio con la stessa causa già accaduto e richiuso in una sessione
precedente. Ricaricato il file corretto e riverificato da zero: ogni numero
riportato qui sotto è stato ricalcolato con il modulo funzionante, non
assunto dalla sessione in cui la regressione era presente.

\verifica{Confronto, a rumore nullo, fra il correlatore dal circuito completo
(preparazione ansatz VQE + Trotter + Hadamard test) e quello del modulo di
riferimento della Parte 1, stessa preparazione, nessun rumore:
$(i,\alpha,j,\beta,t,N)=(2,x,1,x,2.0,20)$ dà $+0.074599+0.257465i$ in entrambi
i casi, scarto $2.5\times10^{-14}$ --- precisione quasi macchina, su questa e
altre due combinazioni indipendenti controllate. L'infedeltà fra ansatz e
stato esatto è $1.6\times10^{-11}$ (dipende dal seme del VQE, non è un
problema); il costo in gate è $18$ CNOT con la preparazione esatta contro $19$
con quella VQE, a $N=5$ --- un solo CNOT in più, trascurabile.}

\section{Il conteggio dei gate cresce come deve}

\begin{center}
\renewcommand{\arraystretch}{1.2}
\begin{tabular}{@{}cc@{}}
\toprule
$N$ & CNOT totali\\
\midrule
$1$ & $7$\\
$2$ & $10$\\
$5$ & $19$\\
$10$ & $34$\\
$20$ & $64$\\
$40$ & $124$\\
\bottomrule
\end{tabular}
\end{center}

Incremento costante di 3 CNOT per unità di $N$ --- il contributo di un passo
di Trotter, lo stesso già visto in Parte 1 e nel Documento 7. Nessuna fusione
spuria: la transpilazione a blocchi funziona come deve anche in questo
circuito più complesso.

\section{Formula di readout: due cammini indipendenti}

\begin{center}
\renewcommand{\arraystretch}{1.2}
\begin{tabular}{@{}lc@{}}
\toprule
Quantità & Valore\\
\midrule
$\langle Z\rangle$ analitico (rumore di gate + correzione readout) & $+0.044089$\\
$\langle Z\rangle$ Monte Carlo ($2\times10^5$ shot, readout vero) & $+0.046550$\\
$|$differenza$|$ (atteso $\lesssim0.0067$, $3\sigma$) & $2.46\times10^{-3}$\\
\bottomrule
\end{tabular}
\end{center}

La formula analitica $(1-2p)\langle Z\rangle_\text{ideale}$, applicata dopo
aver già simulato il rumore di gate, e una simulazione Monte Carlo con un vero
\texttt{ReadoutError} agganciato ad Aer --- due cammini di calcolo indipendenti
--- coincidono entro l'errore statistico atteso.

\section{Risultato: la curva $|C(N)|$ e $N^*$}

\begin{figure}[htbp]
\centering
\includegraphics[width=0.72\textwidth]{figure/fig_passo4_correlatore_vs_N}
\caption{Modulo del correlatore rumoroso $|C_{21}^{xx}(t{=}2,N)|$ in funzione
di $N$. Massimo a $N^*=5$ --- più basso dell'$N^*=8$ della fedeltà di Trotter
(Documento 7). Atteso: il circuito di Hadamard test ha più gate a parità di
$N$, per via dei due blocchi controllati aggiuntivi, quindi il rumore satura
prima.}
\end{figure}

\FloatBarrier

\begin{center}
\renewcommand{\arraystretch}{1.2}
\begin{tabular}{@{}lc@{}}
\toprule
Quantità & Valore\\
\midrule
$N^*$ & $5$\\
$|C(N^*)|$ & $0.250002$\\
\bottomrule
\end{tabular}
\end{center}

\verifica{Rieseguito indipendentemente il correlatore $C_{21}^{xx}(t{=}2)$
per $N=1,5,10,20,40$: a $N=5$, rumore acceso, la parte reale e immaginaria
danno $|C|=\sqrt{0.046662^2+0.245608^2}\approx0.2500$, coerente con il valore
riportato in tabella.}

\messaggio{Tre fatti da questo documento sono usati esplicitamente nel
Documento 9: $N^*=5$ è specifico di questo correlatore, a questi parametri ---
lo scan ne verifica la stabilità, non la dà per scontata; la correzione di
readout è un fattore puramente moltiplicativo, ed è l'argomento (verificato,
non solo dichiarato) per cui $N^*$ non può dipendere da $p_\text{readout}$ per
readout simmetrico; il conteggio gate lineare in $N$ è la base per stimare
come il rumore accumulato cresca al variare dei parametri del canale.}

\end{document}
